\section{Related Work}

\label{sec:related_work}

\subsection{Lexical and Hybrid Retrieval}

Lexical retrieval ranks documents by term frequency and inverse document frequency with length normalization. BM25, the probabilistic lexical ranker used as our primary baseline, has remained competitive with learned dense retrievers across many benchmarks because of its simplicity, robustness to domain shift, and lack of training requirements \citep{ref_9dbf664b6954efbc}. Hybrid retrievers that fuse lexical and dense semantic signals have demonstrated accuracy gains over either signal alone in open-domain question answering and passage retrieval. The agent memory setting differs from classical IR in one structural respect: candidate documents are written by the agent itself during task execution, giving them timestamps and access frequencies that can serve as auxiliary ranking signals alongside lexical overlap. This makes lexical-only baselines like BM25 a particularly informative comparison point: any gap between BM25 and a lexical-plus-temporal scorer measures the net effect of adding time-aware signals.

The position-sensitivity finding of Lost in the Middle provides additional motivation \citep{ref_9dbf664b6954efbc}. In multi-document question answering, relevant information at the middle of a long context is harder for language models to exploit than information at the ends. Retrieval selectors that deliver relevant items first can partially compensate for this bias. CTLS is designed to guarantee that lexically relevant items are packed before irrelevant ones, which aligns with the position-sensitivity finding: if the relevant item must reach the model at all, it should occupy an early slot in the injected context.

\subsection{Agent Memory Architectures}

Memory control and organisation for language-model agents has attracted proposals at increasing levels of autonomy. MemGPT introduces an operating-system-inspired virtual context manager that moves information across memory tiers using explicit management calls, demonstrating document analysis and multi-session chat scenarios \citep{ref_959d6858d034c6e1}. The present study draws on MemGPT's conceptual vocabulary of tiered storage -- a process-local working tier, a session-durable intermediate tier, and a project-persistent long-term tier -- but does not implement its autonomous virtual-memory controller or its interrupt-based flow management.

A-MEM organizes memories as structured notes linked by dynamically computed associations following a Zettelkasten-inspired design \citep{ref_1aacba116cbe3c5b}. Our system uses fixed lexical ranking rather than autonomous linking and evolution, and no empirical comparison to A-MEM was conducted. Generative Agents synthesize natural-language experience records into reflections and dynamically retrieve them for planning behavior in an interactive social simulation \citep{ref_6e953134a3c04402}. We draw on the general retrieval framing but implement neither reflection, planning, nor social simulation.

MemoryBank supports retrieval and updates for sustained personalized interaction, modeling forgetting and reinforcement influenced by elapsed time and importance through an Ebbinghaus-inspired forgetting curve \citep{ref_713e886950406893}. Our temporal-decay term is a simpler exponential heuristic that does not replicate MemoryBank's reinforcement dynamics or its longevity-weighted importance updates. The key architectural difference between these systems and our proposed CTLS is that CTLS makes no assumption about which items were important at write time: it relies entirely on query-time lexical overlap to determine tier membership, which makes the Non-Dominance Property hold unconditionally.

\subsection{Memory Evaluation Benchmarks}

LongMemEval evaluates five core long-term memory abilities -- information extraction, multi-session reasoning, temporal reasoning, knowledge updates, and abstention -- and explicitly separates indexing, retrieval, and reading as distinct evaluation stages \citep{ref_ab2c226739f2f52a}. The stratified subset used in this paper covers a balanced selection from LongMemEval-S with disclosed deviations including a substituted judge model and a losslessly segmented candidate bank; it is not an official benchmark result. The primary comparison is a selector quality comparison on a frozen bank rather than an end-to-end memory system evaluation.

\subsection{Temporal Dynamics in Retrieval}

The interaction between recency and relevance has been studied in information retrieval and recommendation systems. Time-aware retrieval systems often assign higher scores to fresher documents under the assumption that freshness correlates with relevance in news and web search. However, in agent memory settings the assumption fails: conversational facts from months ago are often more relevant to a current query than recent small-talk.

Temporal gating strategies fall on a spectrum from hard gates to soft scoring. Hard temporal gates filter out documents outside a recency window entirely, guaranteeing no stale item enters the ranked list but discarding potentially relevant older material. Soft temporal scoring applies a decay function that reduces but does not eliminate the contribution of older items, preserving recall at the cost of introducing the ordering inversions studied in this paper \citep{ref_temporal_hard_soft_gate}. The hybrid scorer analyzed here is a soft temporal scorer: its decay term $D(m) = 2^{-a(m)/h}$ smoothly downweights older items without excluding them, which is precisely the condition under which temporal dominance can occur.

Recursive retrieval patterns that revisit earlier context to augment or revise retrieved content have been explored in generation tasks and can partially compensate for temporal ordering errors by re-examining candidates whose relevance becomes apparent only after an initial retrieval pass \citep{ref_re3_retrospective_retrieval}. CTLS takes a different approach: rather than compensating after the fact, it prevents the ordering error by construction through the tier partition, ensuring that lexically relevant items are always offered before irrelevant ones regardless of their timestamps.

This mismatch between the freshness heuristic and actual relevance is the core of the temporal dominance failure mode that CTLS addresses. Prior work has not derived a formal condition under which temporal weights cause ranking inversions, nor proposed a tier-partitioned corrective with a provable Non-Dominance Property that holds unconditionally for any weight assignment.

\subsection{Automated Scientific Literature Processing}

Recent work has assembled multi-agent pipelines for autonomous literature review and research paper generation. These pipelines decompose the research task into specialized roles -- retriever, reader, critic, synthesiser -- and pass structured information between agents to produce draft manuscripts from pre-writing materials. Prior systems in this space lack a queryable memory layer with structured provenance across retrieval iterations: once a claim is retrieved and included in a draft, there is no persistent record of which document it came from, when it was retrieved, or how it relates to competing claims. The agent pipeline design described in Section~\ref{sec:method_agent} addresses this gap by equipping each agent with the three-level memory architecture and writing provenance-preserving items to $L_2$. The controlled evaluation in this paper tests the retrieval component of that pipeline in isolation; the end-to-end pipeline evaluation, including the critic's importance updates and the synthesiser's iterative hypothesis refinement, remains future work.


Selecting which retrieved items fit into a fixed prompt budget is a combinatorial problem distinct from ranking quality. Our whole-item first-fit packing policy is a deterministic greedy heuristic. An exact binary knapsack formulation would select the subset of items maximizing total relevance score subject to the token budget constraint; with at most ten candidates the exact solution is tractable in practice. The packing step is held constant across all selectors in our controlled evaluation, so packing differences do not confound the comparison. CTLS affects only the order in which items are offered to the packer, not the packing rule itself.
